\chapter{Eye Pain (Ophthalmalgia)}

\section*{Definition}
Pain or discomfort localized to the eye or orbital region, which may be sharp, dull, burning, or throbbing, originating from ocular or periocular structures.

\section*{What Happens in Your Body}
Eye pain is categorized as ocular (surface: cornea, conjunctiva, sclera) or orbital (deeper: optic nerve, extraocular muscles, vascular structures). Surface pain (foreign body sensation, burning) arises from corneal epithelial disruption (abrasion, dry eye, keratitis) or conjunctival inflammation. Deep, aching pain suggests glaucoma, optic neuritis, uveitis, or orbital pathology. Referred pain from trigeminal nerve branches may mimic eye pain.

\section*{Causes}
\begin{itemize}
  \item Corneal abrasion or foreign body
  \item Dry eye disease (keratoconjunctivitis sicca)
  \item Conjunctivitis (viral, bacterial, allergic)
  \item Keratitis (infectious, contact lens-related)
  \item Acute angle-closure glaucoma (emergency)
  \item Uveitis or iritis
  \item Optic neuritis (associated with multiple sclerosis)
  \item Sinusitis (referred pain to eye)
  \item Orbital cellulitis (emergency)
  \item Cluster headache or migraine (ocular pain)
\end{itemize}

\section*{When to See a Doctor}
See your doctor if:
\begin{itemize}
  \item Symptoms are severe or getting worse over time
  \item Sudden severe eye pain with vision loss, redness, and halos (possible acute glaucoma), eye pain after trauma, or eye pain with fever and proptosis
  \item You have other concerning symptoms such as fever, weight loss, or bleeding
\end{itemize}

\section*{Related Symptoms}
\begin{itemize}
  \item Headache
  \item Vision blurred
  \item Red eye
  \item Photophobia
  \item Nausea
\end{itemize}

\textbf{Important:} If this symptom persists, your doctor will check for serious underlying causes.

\section*{Self-Care / Home Management}
\begin{itemize}
  \item Rest the eyes by taking breaks from screens and near work.
  \item Use artificial tears for surface eye pain related to dryness or irritation.
  \item Apply a cool compress to reduce ocular surface inflammation.
  \item Avoid rubbing the eye, which can worsen corneal abrasions or infections.
  \item Seek immediate medical attention for sudden severe eye pain, vision loss, or eye trauma.
\end{itemize}

\section*{How Your Doctor Will Evaluate You}
\begin{itemize}
  \item Your doctor will measure visual acuity and perform a slit-lamp examination of the anterior segment.
  \item Intraocular pressure is measured to rule out glaucoma.
  \item Pupillary reaction and extraocular movements are assessed to identify neurological causes.
  \item Fluorescein staining is used to detect corneal abrasions, ulcers, or foreign bodies.
  \item CT or MRI of the orbits may be ordered if orbital pathology is suspected.
\end{itemize}

\section*{Treatment Options}
\begin{itemize}
  \item Topical antibiotics for bacterial keratitis or corneal abrasion with infection risk.
  \item Topical corticosteroids for inflammatory conditions (uveitis, scleritis) under specialist supervision.
  \item Immediate treatment for acute angle-closure glaucoma (laser iridotomy, pressure-lowering drops).
  \item Lubricating eye drops and punctal plugs for severe dry eye disease.
  \item Analgesics and abortive therapy for migraine or cluster headache-related eye pain.
\end{itemize}


\section*{References}
\begin{thebibliography}{99}
\bibitem{eye_pain:friedman2009}
Friedman NJ. Eye pain. \textit{Am Fam Physician}. 2009;79(7):579--584.
\bibitem{eye_pain:guluma2012}
Guluma K, Lee JE. Ophthalmology. In: \textit{Rosen's Emergency Medicine}. 8th ed. Elsevier; 2014.
\bibitem{eye_pain:messelink2018}
Messelink MA, van der Sprenkel JW, Vaessen D. The painful eye. \textit{BMJ}. 2018;363:k4406.
\bibitem{eye_pain:kanski2015}
Kanski JJ, Bowling B. \textit{Clinical Ophthalmology: A Systematic Approach}. 8th ed. Elsevier; 2016.
\bibitem{eye_pain:american2019}
American Academy of Ophthalmology. \textit{Preferred Practice Pattern: Primary Open-Angle Glaucoma}. AAO; 2015.
\bibitem{eye_pain:gerstenblith2012}
Gerstenblith AT, Rabinowitz MP. \textit{The Wills Eye Manual}. 7th ed. Lippincott Williams \& Wilkins; 2016.
\end{thebibliography}
